% ============================================================
%  The Lagrangian OSSE --- design document
%  Standalone: compile with pdflatex.
%  Packages used: amsmath, amssymb, amsthm, booktabs, geometry,
%                 hyperref, enumitem, longtable
% ============================================================

\documentclass[11pt]{article}

\usepackage[margin=1in]{geometry}
\usepackage{amsmath,amssymb,amsthm}
\usepackage{booktabs}
\usepackage{enumitem}
\usepackage{longtable}
\usepackage[hidelinks]{hyperref}

\newtheorem{remark}{Remark}
\newtheorem{assumption}{Assumption}

\title{The Lagrangian OSSE\\[2mm]
\large Measuring preferential sampling bias in regional ocean heat content\\
by simulating floats inside a model ocean}
\author{}
\date{}

\begin{document}
\maketitle

% Long CMEMS dataset identifiers in \texttt do not hyphenate; relax line breaking.
\sloppy

%=============================================================
\section{Why this experiment exists}
%=============================================================

Preferential sampling bias is, in the real ocean, unobservable by construction: it is the
difference between what floats measured and what they would have measured had they gone
somewhere else. It cannot be estimated from the profiles themselves
(Section~\ref{sec:notestimable}), so every observational approach to it is indirect.

The Observing System Simulation Experiment (OSSE) removes that obstruction by replacing the
ocean with a model in which the field is known everywhere. Synthetic floats are advected
through the model's own velocity field, sample the model's own heat content, and the bias is
computed by subtraction rather than inferred. Three things follow that are unavailable
otherwise:

\begin{enumerate}[leftmargin=*]
  \item \textbf{Magnitude.} A number for the bias, with an ensemble error bar.
  \item \textbf{Validation.} The observational estimator---eddy-occupancy
        post-stratification---can be applied to the synthetic floats and checked against the
        known answer. Without this, that estimator is untested.
  \item \textbf{Mechanism.} Direct association between float behaviour (residence time,
        looping) and the anomaly encountered, measurable per float.
\end{enumerate}

The cost is that the answer is a statement about the model. Section~\ref{sec:transfer}
addresses what transfers and what does not.

%=============================================================
\section{Theoretical background}
%=============================================================

\subsection{Target, estimator, error}

Let $\mathcal{R} = D \times \mathcal{T}$ be the analysis region: $D$ a horizontal domain masked
to bathymetry deeper than $2000$\,m, $\mathcal{T}$ an averaging window. Let $P$ be normalised
uniform (area $\times$ time) measure and $Y : \mathcal{R} \to \mathbb{R}$ the $0$--$2000$\,m
ocean heat content (OHC) in J\,m$^{-2}$. Profiles arrive at $z_1,\dots,z_n$, with empirical
sampling measure $Q_n = n^{-1}\sum_j \delta_{z_j}$. Then
\begin{equation}
  \bar{Y} = \int_{\mathcal{R}} Y \, dP,
  \qquad
  \hat{Y} = \int_{\mathcal{R}} Y \, dQ_n,
  \qquad
  \hat{Y} - \bar{Y} = \int_{\mathcal{R}} Y \, d(Q_n - P).
\end{equation}

\subsection{The modelling lineage}
\label{sec:lineage}

The decomposition adopted below is not new. It is the standard on which the current generation
of Argo-based OHC products rests, and setting out the lineage matters because it fixes the
scope of the present work: the assumption under examination is common to all of these products
and is not an idiosyncrasy of any one of them.

\begin{description}[leftmargin=*,style=nextline]
  \item[Roemmich \& Gilson (2009).] Objective mapping of Argo temperature and salinity onto a
    regular grid, producing the widely used RG climatology and anomaly product. Establishes the
    working practice of a climatological mean field plus an interpolated anomaly.
  \item[Kuusela \& Stein (2018).] Replaces global stationarity with a locally stationary
    spatio-temporal Gaussian process, with covariance parameters estimated in moving windows.
    This is the mapping foundation on which the subsequent OHC work explicitly builds.
  \item[Park, Kuusela, Giglio \& Gray (2023).] Applies a local latent Gaussian process
    framework to ocean heat \emph{transport}, introduces an approximate EM algorithm estimating
    the mean field and covariance parameters jointly, and adds a debiasing step for the mean
    field. A sibling application rather than a direct antecedent for heat \emph{content}, but it
    is where joint mean/covariance estimation enters the lineage.
  \item[Sukianto et al.\ (2026).] An end-to-end mapping and uncertainty quantification framework
    for OHC built on the Kuusela--Stein foundation, with data-driven decorrelation scales and
    variances and uncertainty quantification by local conditional simulation ensembles.
    Underlies the LocalGP data product, which has contributed to climate assessments and OHC
    intercomparison efforts.
  \item[Bivariate extension (2026).] Joint modelling of two pressure layers, allowing the more
    sparsely sampled midocean layer to borrow predictive strength from the upper ocean.
\end{description}

Three remarks about this lineage bear directly on the design of the present experiment.

\begin{remark}[``Debiased'' in Park et al.\ refers to mean-field misspecification]
\label{rem:debias}
The debiasing procedure of Park et al.\ refines a potentially under-specified parametric mean
field, correcting leakage of mean structure into the estimated anomaly. In the terminology of
Section~\ref{sec:coverage-pref} this is a refinement on the \emph{coverage} side: it improves
the description of structure that is a deterministic function of position. It does not address
dependence between the sampling design and the field, and the two should not be conflated.
\end{remark}

\begin{remark}[The mean field is Argo-derived]
\label{rem:circular}
Both the EM procedure of Park et al.\ and the plug-in estimates of Sukianto et al.\ estimate
$\mu$ from the same profiles being mapped. If the array preferentially samples warm water,
$\hat\mu$ inherits that bias. Preferential sampling is therefore not confined to the nugget at
sample points: it propagates into the reference field against which anomalies are defined. In
the real analysis this argues for an Argo-independent climatology; inside the OSSE, where the
model supplies its own truth, the issue is handled by internal closure
(Section~\ref{sec:consistency}).
\end{remark}

\begin{remark}[What is actually at risk]
\label{rem:trend}
A preferential offset that is roughly constant in time affects the estimated mean state but
largely cancels in anomalies and trends---the quantities used in Earth energy imbalance work.
The threat to trends is \emph{time-varying} preferentiality. The array roughly tripled between
2004 and 2007; parking depths were standardised; float models and mission parameters changed;
and Southern Ocean coverage expanded substantially. If trapping propensity or eddy-class
occupancy imbalance drifted over the record, the consequence is a spurious OHC trend. This
motivates a direct diagnostic---compute the occupancy imbalance year by year and test for
trend---which is inexpensive once the eddy census exists and which we include among the
deliverables (Section~\ref{sec:deliverables}).
\end{remark}

\subsection{The standard statistical model}
\label{sec:model}

We adopt without modification the decomposition used in that literature:
\begin{equation}
  \label{eq:model}
  Y(x,y,t) \;=\; \mu(x,y,t) \;+\; a(x,y,t) \;+\; \varepsilon(x,y,t),
\end{equation}
where
\begin{itemize}[leftmargin=*]
  \item $\mu$ is a climatological mean field, estimated by local polynomial regression on a
        regular grid;
  \item $a$ is a zero-mean transient anomaly field, modelled as a locally stationary
        spatio-temporal Gaussian process;
  \item $\varepsilon$ is a zero-mean nugget describing unresolved fine-scale variation,
        principally mesoscale, taken as uncorrelated in space and time within the local window.
\end{itemize}

Nothing below changes these definitions, the estimation of $\mu$, or the locally stationary
Gaussian process for $a$. The mapped products built on \eqref{eq:model} are taken as given.

\subsection{Preferential sampling is compatible with this model}
\label{sec:compatible}

It is worth being explicit about why, since \eqref{eq:model} is often read as excluding the
phenomenon.

\paragraph{``Zero-mean'' is a marginal statement at a fixed location.}
$\mathbb{E}[\varepsilon(z)] = 0$ says nothing about $\mathbb{E}[\varepsilon(z_j) \mid
\mathcal{Z}]$, a conditional expectation at a \emph{random, field-dependent} location. Toy
case: let $\varepsilon$ be i.i.d.\ $N(0,\sigma^2)$ on a grid and let the sampler select the
largest value. The marginal mean is zero at every grid point; the mean at the sampled point is
positive. No assumption in \eqref{eq:model} is violated.

\paragraph{``Uncorrelated'' is likewise a statement about pairs at fixed locations.}
$\operatorname{Cov}(\varepsilon(z),\varepsilon(z')) = 0$ constrains the joint law of the field
at prescribed points and places no constraint on which points get sampled.

\paragraph{What is added is one conditional distribution.}
The standard framework implicitly factorises the joint law as
\[
  [\,a,\varepsilon\,] \;\times\; [\,\mathcal{Z}\,] \;\times\;
  [\,\text{data} \mid \mathcal{Z}, a, \varepsilon\,],
\]
with the design $\mathcal{Z}$ free-standing. We replace $[\mathcal{Z}]$ by
$[\mathcal{Z} \mid a, \varepsilon]$. That is the entire modification.

\paragraph{Why the modification matters.}
The target and the estimator treat $\varepsilon$ asymmetrically. In the target
$\int (\mu + a)\,dP$, the nugget integrates away over a region spanning many correlation
lengths. In the estimator $n^{-1}\sum_j [\mu(z_j) + a(z_j) + \varepsilon(z_j)]$ it does
\emph{not} average out, if the $z_j$ were selected on it. So $\varepsilon$ contributes nothing
to what is wanted and something to what is computed. Write
\begin{equation}
  \label{eq:b}
  b_j \;:=\; \mathbb{E}\!\left[\varepsilon(z_j) \mid \mathcal{Z}\right]
\end{equation}
for the conditional mean nugget at the realised design. The standard framework assumes
$b \equiv 0$; this experiment measures it.

\paragraph{A physical tension, and its resolution.}
Formally, white noise and preferential sampling coexist without difficulty. Physically they are
in tension: a float follows a continuous trajectory and cannot steer toward a
delta-correlated field, since there is nothing coherent to be trapped by. The mechanism
requires $\varepsilon$ to have genuine spatial structure.

The resolution is that ``uncorrelated'' in \eqref{eq:model} means uncorrelated \emph{at the
resolution of the mapping}, not literally zero correlation length. Unresolved mesoscale
variability has $\ell_\varepsilon \sim 50$--$100$\,km and a few weeks---negligible against the
data-driven decorrelation scales used for $a$, but far from zero, and precisely the scale at
which eddy trapping operates. We therefore treat $\varepsilon$ as white for mapping purposes
and as short-range correlated with $\ell_\varepsilon > 0$ for the sampling analysis. The same
fact drives the effective-sample-size problem in Section~\ref{sec:control}.

\paragraph{Dependence is restricted to $\varepsilon$, not $a$.}
We assume the design may depend on the unresolved field but not on the resolved anomaly. This
is physically defensible---trapping is a mesoscale phenomenon, mesoscale variability is in
$\varepsilon$ by construction, and large-scale anomalies do not trap floats---and it is
analytically valuable, because it leaves the mapped $a$ usable as given. Were the design to
depend on $a$, a Gaussian process predictor for $a$ would itself be biased, since it conditions
on the data as though locations were exogenous (Diggle et al.\ 2010; Gelfand et al.\ 2012). The
second-order caveat is that a warm large-scale state may coincide with a stronger jet and
faster flushing; we note it and do not model it.

\subsection{Coverage versus preferential}
\label{sec:coverage-pref}

Since the region spans many correlation lengths of $\varepsilon$, we have
$\int_{\mathcal{R}} \varepsilon \, dP \approx 0$, and
\begin{equation}
  \label{eq:split}
  \hat{Y} - \bar{Y}
  \;=\; \underbrace{\int_{\mathcal{R}} (\mu + a)\, d(Q_n - P)}_{\textbf{coverage}}
  \;+\; \underbrace{\frac{1}{n}\sum_{j=1}^{n} \varepsilon(z_j)}_{\textbf{preferential}} .
\end{equation}

\emph{Coverage bias} follows from floats being unevenly distributed in space and time. It
depends on the field only through structure that is a known or estimable function of position,
so given the position list it is computable. It would be present even if float positions were
fixed in advance, independently of the ocean.

\emph{Preferential bias} is the mean of the nugget at the sampled points. It is nonzero only
when the positions are stochastically dependent on the field---floats lingering where water is
anomalously warm, or being flushed from regions where it is anomalously cold. It survives even
under perfectly uniform sampling density.

The mechanism making this plausible for Argo: floats drift with ocean currents, and currents
are geostrophically tied through thermal wind to the density structure that sets heat content.
The instrument is transported by the field it measures.

\subsection{A note on exactness}
\label{sec:exactness}

Equation \eqref{eq:split} is exact if the split between $a$ and $\varepsilon$ is made a
definition---that is, if $a := \Pi_L(Y - \mu)$ and $\varepsilon := (I - \Pi_L)(Y - \mu)$ for
$\Pi_L$ the $L^2(P)$-orthogonal projection onto a resolved space $V_L$ containing the
constants, in which case $\int \varepsilon \, dP = 0$ identically. In the OSSE we adopt this
construction, since the model field is available everywhere and the projection is computable.
In the operational setting the split is instead induced by the fitted Gaussian process, and
\eqref{eq:split} holds approximately, with $\int \varepsilon \, dP$ small but not exactly zero.
Reporting the OSSE's projection residual quantifies that approximation.

\subsection{Why the preferential term is not estimable from profiles}
\label{sec:notestimable}

Two obstructions:

\begin{itemize}[leftmargin=*]
  \item The residuals available in practice, $\hat\varepsilon_j$, come from fitting $\mu + a$
        to the observed profiles. When the constant lies in the fitted span,
        $\sum_j \hat\varepsilon_j = 0$ identically.
  \item More fundamentally, $\varepsilon$ is observed only where floats went. If floats
        systematically avoided cold water, no feature of the sample reveals the temperature of
        the water avoided.
\end{itemize}

Note also that \textbf{cross-validation cannot detect this}. Holding out observed profiles and
predicting them from other observed profiles validates against a sample carrying the same
selection as the training set. A map can pass cross-validation cleanly and remain biased.

Hence the OSSE: measure $b$ where $\varepsilon$ is known everywhere.

\subsection{Why the OSSE targets $\varepsilon$, not raw $Y$}
\label{sec:whyresidual}

\paragraph{Recoverability.} $\mu$ is a known function of position, so its contribution to
\eqref{eq:split} is computable from the position list alone, and $a$ is estimable from the
data. Only $\varepsilon$ leaves a residue that nothing recovers. Raw OHC bundles all three and
measures something largely already correctable.

\paragraph{Signal-to-noise.} Across an energetic box, $\mu$ varies by $O(10^{9}$--$10^{10})$
J\,m$^{-2}$ while the preferential target is $O(10^{8})$. A statistic built on raw OHC is
dominated by mean-field sampling and demands far more ensemble members to resolve the term of
interest.

\paragraph{Nonetheless, report the total too.} In the OSSE the truth is known, so the full
error $\hat{Y} - \bar{Y}$ is free. Its value: the total is invariant to where the
resolved/unresolved split is drawn, whereas the coverage/preferential split is not. Enlarging
the resolved space moves variance from $\varepsilon$ into $a$ and shrinks the preferential
term. Reporting both makes that dependence visible rather than hidden.

\paragraph{What the nugget contains.} Argo temperature accuracy of $\sim 0.002\,^\circ$C
integrated over $2000$\,m corresponds to $\sim 1.6\times 10^{7}$ J\,m$^{-2}$, against mesoscale
OHC variability of order $10^{9}$. Instrument error is thus a percent or two of the nugget, and
it is genuinely design-independent---a CTD's accuracy does not depend on which eddy it sits in.
The remaining $\sim\!98\%$ is geophysical, and it is exactly the component for which the
independence assumption is least secure.

%=============================================================
\section{Design}
%=============================================================

\subsection{Core principle}

Inside the OSSE, \textbf{the model is the ocean}. Every quantity in
Section~\ref{sec:model} is computed from model fields, with no reference to the real ocean or
to real observations. This closure is what makes the experiment interpretable; violations of it
(Section~\ref{sec:consistency}) are the main way it goes wrong.
\begin{align}
  \bar{Y}^{\text{sim}} &= \int_{\mathcal{R}} Y^{\text{model}} \, dP
    &&\text{(exact, by quadrature over the model grid)} \\
  \hat{Y}^{\text{sim}} &= \frac{1}{n}\sum_j Y^{\text{model}}(z_j^{\text{sim}}) && \\
  b^{\text{sim}} &= \frac{1}{n}\sum_j \varepsilon^{\text{model}}(z_j^{\text{sim}}) &&
\end{align}
with $\varepsilon^{\text{model}} = (I - \Pi_L)\big(Y^{\text{model}} - \mu^{\text{model}}\big)$.

\subsection{Model choice}
\label{sec:modelchoice}

\begin{table}[h]
\centering
\begin{tabular}{lll}
\toprule
 & GLORYS12 & LLC4320 \\
\midrule
Resolution & $1/12^\circ$ ($\sim$6 km at $40^\circ$N) & $1/48^\circ$ ($\sim$1.5 km) \\
Period & 1993--present & Sep 2011 -- Nov 2012 \\
Assimilation & Yes (incl.\ Argo) & Free-running \\
Lagrangian realism & \emph{damped} & honest \\
Period matching & can match 2015--2019 & cannot \\
\bottomrule
\end{tabular}
\caption{Candidate models. Assimilation increments disrupt eddy coherence and shorten
Lagrangian trapping times, so the reanalysis is expected to understate the effect.}
\end{table}

\textbf{Use both, for different purposes.} The free-running model gives trustworthy Lagrangian
statistics and is the primary estimate of the mechanism. The assimilating reanalysis gives a
period-matched sensitivity run. Expect GLORYS to \emph{understate} $b^{\text{sim}}$; report the
pair and treat the gap as part of the uncertainty.

\textbf{A separate use to avoid.} A reanalysis that assimilates Argo must not be used for a
resampling test on \emph{real} float positions: at those positions the field is partly the
float's own measurement propagated by the assimilation covariance, so shifting the trajectory
removes a constraint and produces a difference with no oceanographic content. That confound
does not arise here, because the OSSE uses only synthetic floats and the model's own truth.

\subsection{Float simulation}
\label{sec:simulation}

Integrate parcel trajectories through the model's three-dimensional velocity, reproducing the
Argo park-and-profile cycle.

\begin{table}[h]
\centering
\begin{tabular}{llll}
\toprule
Phase & Duration & Depth & Advecting velocity \\
\midrule
Park & $\sim$9 days & 1000 m & $\mathbf{u}(1000\,\mathrm{m})$ \\
Descend & $\sim$6 h & 1000 $\to$ 2000 m & interpolated through the column \\
Profile ascent & $\sim$6 h & 2000 m $\to$ surface & sheared flow through the column \\
Surface drift & 6--24 h & 0 m & surface current (plus optional windage) \\
Descend & $\sim$6 h & 0 $\to$ 1000 m & interpolated \\
\bottomrule
\end{tabular}
\caption{Simulated float cycle.}
\end{table}

Three details matter:

\begin{itemize}[leftmargin=*]
  \item \textbf{Do not treat the profile as instantaneous.} Ascent from $2000$\,m at
        $\sim$10\,cm\,s$^{-1}$ takes $\sim$6\,h; in a strong current the float moves tens of km
        through vertically sheared flow. The profile's nominal position is not the parking
        position.
  \item \textbf{Surface drift is not negligible.} Surface velocities typically exceed
        $1000$\,m velocities substantially, and the surface interval is where wind enters.
  \item \textbf{The box is not closed.} Floats advect in and out. Seed over a wider domain and
        count profiles falling inside $D$.
\end{itemize}

\paragraph{Implementation.} Use an established Lagrangian framework rather than writing an
integrator: OceanParcels (Python; handles NEMO-family output, which includes GLORYS, natively),
Ariane, or TRACMASS. Fourth-order Runge--Kutta with sub-daily timestep, trilinear spatial and
linear temporal interpolation of velocity.

\paragraph{Seeding and mortality.} For each real float present in the box during the period,
seed a synthetic float at the corresponding deployment position and date. Impose a per-cycle
failure probability tuned so that simulated profile counts match the real record. Target
$n \approx 8500$ profiles over five years in a $4.3\times10^{6}$\,km$^2$ box, from $\sim$45--50
resident floats.

\paragraph{Ensemble.} Generate members by perturbing seed positions (a few km) and dates (a few
days), and by varying the mortality seed. Sizing: with per-realisation noise
$\sigma/\sqrt{n_{\text{eff}}} \approx 6\times10^{7}$ J\,m$^{-2}$ (taking $\sigma \sim 10^{9}$,
$n_{\text{eff}} \sim 300$), resolving $b \sim 5\times10^{7}$ at three standard errors needs
roughly 12 members; \textbf{use 30--50}, pooling across years where the model period permits.

\subsection{The decoupled control}
\label{sec:control}

$b^{\text{sim}}$ as computed contains both terms of \eqref{eq:split}. To isolate the
preferential component, build a control in which trajectories are kinematically realistic but
carry \textbf{no alignment with the field being sampled}.

\paragraph{Cheap construction (recommended).} Reuse the same trajectories, but evaluate
$\varepsilon^{\text{model}}$ from a \emph{different year} of the model at the same day-of-year
and position:
\begin{equation}
  b^{\text{ctrl}} = \frac{1}{n}\sum_j \varepsilon^{\text{model}}_{(\text{year } y')}(z_j^{\text{sim}}),
  \qquad y' \neq y .
\end{equation}
Float positions were determined by year $y$'s eddies, so with respect to year $y'$'s field they
are exogenous. This costs nothing extra---no re-integration---and averaging over several $y'$
gives the control's own sampling distribution directly. Then
\begin{equation}
  \widehat{\mathcal{P}}^{\text{sim}} = b^{\text{sim}} - \overline{b^{\text{ctrl}}},
  \qquad
  \widehat{\mathcal{C}}^{\text{sim}} = \overline{b^{\text{ctrl}}},
\end{equation}
and the spread of $b^{\text{ctrl}}$ across $y'$ is the null distribution against which
$b^{\text{sim}}$ is tested.

\paragraph{Why trajectories are preserved rather than resampled.} The null distribution must
answer: \emph{if sampling were exogenous, how much would the statistic vary?} Its width is set
by the number of genuinely independent looks at the field. Because a float moves only
$30$--$80$\,km per 10-day cycle---well inside a mesoscale feature---consecutive profiles are
near-replicates and $n_{\text{eff}}$ is far below $n$. Drawing i.i.d.\ points from an estimated
intensity would give a null several times too narrow and would reject on a perfectly exogenous
array. Reusing the actual trajectories reproduces the correct width automatically, so
$n_{\text{eff}}$ never has to be estimated. It also delivers the correct non-Gaussian shape of
the null, which matters at a few hundred effective looks.

This preserves \emph{kinematic} persistence (bounded float speed, smooth tracks) while breaking
\emph{field alignment} (loops sitting on warm eddies). Only the second is the mechanism under
test. Conditioning on trajectory shape does cost power, since shape is partly caused by the
field---but it can only make the test conservative, not produce a false positive.

\paragraph{Alternative.} Advect a second ensemble using year $y'$ velocity while sampling year
$y$ field. More expensive, and preferable only if the control's trajectory geometry must also
be independent.

\subsection{Day-of-year alignment}

Use year-shifts rather than circular time-shifts within a single window. Seasonality is matched
automatically, replicates are genuinely independent rather than wrap-around artifacts, and a
long model record yields many more of them.

%=============================================================
\section{Analysis domain and data}
\label{sec:data}
%=============================================================

\subsection{Period: 2015--2019}
\label{sec:period}

\begin{table}[h]
\centering
\begin{tabular}{p{0.28\textwidth}p{0.62\textwidth}}
\toprule
Criterion & Rationale \\
\midrule
Mature array & Deployment reached $\sim$3800 floats by 2007; count and spatial coverage are
stable across the window, so occupancy statistics are not confounded by array growth. \\
Pre-disruption & Deployments were cut sharply from 2020 and recovery was uneven by basin. \\
Delayed-mode complete & Quality-controlled profiles are fully available, unlike the most
recent one to two years. \\
Sufficient sampling & $\sim$1700 profiles per box per year, $\sim$8500 over five years, from
$\sim$45--50 resident floats. \\
Averaging helps twice & The five-year mean field is smoother than any single year, and $n$ is
five times larger; both shrink the coverage term, so a null result on the preferential term is
more meaningful. \\
\bottomrule
\end{tabular}
\caption{Rationale for the 2015--2019 analysis window.}
\end{table}

A five-year window is preferred to a single year because the single-year case is both harder
and more likely to be uninformative; it is the better second experiment, not the first.

\paragraph{A scheduling conflict, and its resolution.} LLC4320 covers only September 2011 to
November 2012, so the free-running arm of Section~\ref{sec:modelchoice} cannot be
period-matched to 2015--2019. This is acceptable because the OSSE's purpose is to validate the
eddy-occupancy estimator \eqref{eq:poststrat}, and a validated estimator transfers across
periods. The arrangement should be stated explicitly: the free-running run sits at its own
period, GLORYS supplies the period-matched sensitivity run, and the real-array analysis sits at
2015--2019.

\paragraph{A longer record for one diagnostic.} Remark~\ref{rem:trend} concerns
\emph{time-varying} preferentiality, which five years cannot resolve. The occupancy census
requires only float positions and an altimetric eddy atlas---no model, no trajectories---so it
should be run separately over the full Argo record, 2004--2024. A trend in occupancy imbalance
would be a stronger result than the five-year magnitude, and it does not depend on the OSSE.

\subsection{Region: two paired North Atlantic boxes}
\label{sec:region}

\paragraph{Why the North Atlantic.} It is the best-sampled basin, so the array is close to its
design density and one is not measuring a coverage pathology; altimetric coverage for the eddy
census is excellent; the front is strong and its eddy coherence well characterised; and there
is no sea ice. Ice creates regions of zero float coverage, where the framework degenerates
entirely---a real failure mode, but one belonging to the global extension rather than the
pilot.

\paragraph{Why two boxes.} The achievable bias scales with the variability of the field being
selected on, and mesoscale variance differs by roughly a factor of three between quiescent and
frontal regions. A single number is therefore hard to interpret; the \emph{ratio} is the
generalisable result, and it comes from running the same pipeline twice.

\begin{table}[h]
\centering
\begin{tabular}{p{0.14\textwidth}p{0.38\textwidth}p{0.38\textwidth}}
\toprule
 & Box A --- Sargasso & Box B --- Gulf Stream extension \\
\midrule
Nominal extent & 22--34$^\circ$N, 40--60$^\circ$W & 35--45$^\circ$N, 40--60$^\circ$W \\
Character & abyssal plain, quiescent, uniformly deeper than 2000\,m & strong front, high eddy
kinetic energy, rapid advection \\
Mechanism & floats meander quasi-diffusively & floats trapped in rings or flushed downstream
within weeks \\
Expected role & small effect; establishes that the method works & large effect; the substantive
finding \\
\bottomrule
\end{tabular}
\caption{Paired analysis regions, each $\approx 4.3\times10^{6}$\,km$^2$.}
\end{table}

Box A functions as a control: if the pipeline reports a large effect there, something is wrong.

\paragraph{Bathymetry sets the edges---inspect the mask before committing.} These longitudes
lie west of the Mid-Atlantic Ridge (crest near $29$--$30^\circ$W at these latitudes), which is
why they were chosen: the ridge runs $2500$--$3000$\,m with shallower peaks and would punch
holes in the $>2000$\,m mask. However, the New England Seamounts extend southeast to roughly
$35^\circ$N, $57^\circ$W with summits well above $2000$\,m, so the southwestern corner of Box B
may clip them. Either trim the western boundary to $\sim55^\circ$W or mask and report the
excluded area fraction. Equal box areas matter more than the exact coordinates.

\paragraph{Scale.} Each box is large enough to span many correlation lengths of $\varepsilon$,
so that $\int_{\mathcal{R}} \varepsilon \, dP \approx 0$ holds as assumed in
\eqref{eq:split}, and small enough that every quantity can be cross-checked by brute force.

\subsection{GLORYS12 datasets and variables}
\label{sec:glorys}

The reanalysis arm draws on GLOBAL\_MULTIYEAR\_PHY\_001\_030 (GLORYS12V1), a $1/12^\circ$,
50-level global product covering 1993 to the present, distributed as four datasets. Different
parts of the protocol require different ones.

\paragraph{Static fields --- download once, first.}
\texttt{cmems\_mod\_glo\_phy\_my\_0.083deg\_static}

\begin{table}[h]
\centering
\begin{tabular}{ll}
\toprule
Variable & Purpose \\
\midrule
\texttt{deptho} & bathymetry; defines the $>2000$\,m mask \\
\texttt{e3t} & cell thickness; \textbf{essential} for $\int \rho c_p T\,dz$ \\
\texttt{e1t}, \texttt{e2t} & cell dimensions for the area-weighted quadrature of $\int Y\,dP$ \\
\texttt{mask} & land--sea mask \\
\texttt{deptho\_lev} & model level at sea floor; partial cells \\
\bottomrule
\end{tabular}
\end{table}

Do not reconstruct layer thicknesses from the depth coordinate; use \texttt{e3t}. (The Product
User Manual lists its units as m\,s$^{-1}$, which is a typographical error---the standard name
is \texttt{cell\_thickness}, in metres.)

\paragraph{Monthly means --- the source for $\mu$.}
\texttt{cmems\_mod\_glo\_phy\_my\_0.083deg\_P1M-m}, variables \texttt{thetao} (levels spanning
$0$--$2000$\,m) and \texttt{zos}, from 1993 through the end of the analysis window, subset to
the two boxes plus a halo.

This, rather than the packaged climatology, is the right source. The climatology product
supplies a seasonal cycle and long-term mean but \emph{no trend}, and per
Remark~\ref{rem:trend} the trend is precisely what is at risk. Fitting $\mu$ as mean plus trend
plus annual and semiannual harmonics per cell, from monthly means, is also what
Section~\ref{sec:consistency} requires---that $\mu$ be computed the way the operational one is.
At roughly 324 monthly time steps over two boxes this is inexpensive.

The halo need only be a couple of regression bandwidths wide, since $\mu$ is evaluated only at
profiles inside $D$ and over $D$ for the quadrature. It is not the trajectory buffer.

\texttt{zos} at monthly resolution over the same span is two-dimensional and nearly free, and
supplies a long-baseline sea surface height mean for the eddy census---preferable to a mean
computed from the five years of daily fields alone.

\paragraph{Daily means --- the workhorse.}
\texttt{cmems\_mod\_glo\_phy\_my\_0.083deg\_P1D-m}, over the analysis window and a buffered
region.

\begin{table}[h]
\centering
\begin{tabular}{lll}
\toprule
Variable & Levels & Purpose \\
\midrule
\texttt{thetao} & $0$--$2000$\,m & the OHC field \\
\texttt{uo}, \texttt{vo} & restricted; see below & trajectory integration \\
\texttt{zos} & 2D & eddy census \\
\texttt{mlotst} & 2D & optional context \\
\bottomrule
\end{tabular}
\end{table}

Omit \texttt{bottomT}, \texttt{siconc}, \texttt{sithick}, \texttt{usi}, \texttt{vsi}: there is
no ice in either box and sea floor temperature is not used.

\paragraph{Velocity subsetting is the dominant cost decision.} Full-column daily
\texttt{uo}/\texttt{vo} exceeds everything else combined. For a first pass, retain only the
levels bracketing the $1000$\,m parking depth together with the top two or three for the
surface drift---roughly a $90\%$ reduction. The nine-day park dominates displacement; ascent
and descent shear is a refinement. Extend to the full column only if the association in
Figure~5 indicates the mechanism is live.

\paragraph{Salinity.} Omit \texttt{so} initially and use constant $\rho c_p$, as most
operational OHC products do; including it doubles the three-dimensional volume. Download a
single year to quantify the in-situ density effect (of order $1\%$) and then decide.

\paragraph{Depth levels.} The $2000$\,m horizon falls between approximately the fortieth and
forty-first of the fifty levels, so roughly the top forty are required. Verify against the
actual \texttt{depth} coordinate and \texttt{e3t} rather than assuming, and partial-weight the
cell straddling $2000$\,m---at these depths that cell is $\sim275$\,m thick, so mishandling it
is a material error.

\paragraph{Climatology product.} \texttt{cmems\_mod\_glo\_phy\_my\_0.083deg-climatology\_P1M-m}
provides the 1993--2016 monthly mean fields. Useful as a first-pass shortcut and cross-check
only. Its baseline differs from the fitted $\mu$, so expect a constant offset: compare
structure, not absolute values.

\paragraph{Volume.} Global daily files are $3.4$\,GB each, so five years of unsubset data is of
order $6$\,TB. A buffered region of $15$--$50^\circ$N, $30$--$70^\circ$W covers both boxes with
room for floats advecting in and out and is roughly $0.3\%$ of global area. For scale, one box
at $20^\circ\times12^\circ$ with forty levels and $\sim$324 monthly \texttt{thetao} steps is
about $1$\,GB stored as \texttt{short}; both boxes with halo is a few GB, trivial beside the
daily velocity.

\subsection{Practical notes}
\label{sec:practical}

\begin{description}[leftmargin=*,style=nextline]
  \item[The distributed product is not a C-grid.] Fields are interpolated from the native
    Arakawa C-grid onto a collocated regular $1/12^\circ$ grid, with all variables at the same
    points. Trajectory integration must therefore use A-grid interpolation. Applying C-grid
    interpolation to collocated data is a silent error that manifests as spurious divergence in
    the trajectories.
  \item[Packed integers.] All fields are stored as \texttt{short} with \texttt{scale\_factor}
    and \texttt{add\_offset}, and \texttt{\_FillValue} $=-32767$ over land. Standard NetCDF
    readers decode automatically; anything reading raw must apply
    $\text{real} = \text{stored} \times \text{scale} + \text{offset}$.
  \item[Reanalysis spin-up.] GLORYS12 is initialised in January 1993 from temperature and
    salinity regressed on EN4.2.0, and reanalyses commonly drift in deep heat content while the
    assimilation settles. That drift enters the fitted trend coefficient directly. Plot box-mean
    $0$--$2000$\,m OHC against time before fitting, and consider beginning the fit in 1995;
    nonlinearity in the earliest years is a spin-up artifact rather than ocean signal.
  \item[Linearity of the vertical integral.] With constant $\rho c_p$, OHC is linear in
    temperature, so monthly-mean OHC equals OHC computed from monthly-mean temperature. Fitting
    $\mu$ from monthly files therefore introduces no error relative to fitting from dailies.
    This breaks if in-situ density is used---a further small argument for constant $\rho c_p$.
  \item[Baseline choice moves the split, not the total.] Fitting $\mu$ over a window that
    includes the analysis period is not circular here, since $\mu$ derives from the complete
    model field rather than from sampled profiles, so no preferential bias can enter (contrast
    Remark~\ref{rem:circular}). But where the trend is absorbed shifts the $\mu$/$a$ boundary
    and hence the coverage/preferential split; only the total is invariant. Match whatever the
    operational pipeline does and record the choice.
  \item[Access.] Use the \texttt{copernicusmarine} Python toolbox; \texttt{motuclient} is
    retired. A single \texttt{copernicusmarine subset} call accepts variable list, depth range,
    bounding box, and time range. Lazy Zarr access via \texttt{open\_dataset} is worth using for
    the monthly $\mu$ fit even where the trajectory code requires local files.
\end{description}

%=============================================================
\section{Estimands}
%=============================================================

\begin{table}[h]
\centering
\begin{tabular}{lll}
\toprule
Symbol & Definition & Read from \\
\midrule
$b^{\text{sim}}$ & mean nugget at synthetic profile positions & main ensemble \\
$\widehat{\mathcal{C}}^{\text{sim}}$ & coverage component & control mean \\
$\widehat{\mathcal{P}}^{\text{sim}}$ & preferential component & main minus control \\
$\hat{Y} - \bar{Y}$ & total error, all components & main ensemble, raw OHC \\
$\hat{f}_\kappa$ vs $f_\kappa$ & eddy-class occupancy vs areal fraction & Section~\ref{sec:validation} \\
$n_{\text{eff}}$ & effective sample size & $\sigma^2/n$ over control variance \\
$\ell_\varepsilon$ & nugget correlation range & residual variograms \\
\bottomrule
\end{tabular}
\end{table}

\paragraph{Budget closure.} Equation \eqref{eq:split} is an identity under the projection
construction of Section~\ref{sec:exactness}, so
\[
  (\hat{Y} - \bar{Y}) - \big[\widehat{\mathcal{C}}^{\text{sim}} + \widehat{\mathcal{P}}^{\text{sim}}\big] = 0
\]
up to quadrature and interpolation error. A material gap indicates an inconsistency in how
$\Pi_L$ was applied, or a resolved space too small for the coverage term to be captured. Report
the closure residual as a numerical check.

%=============================================================
\section{The validation payoff}
\label{sec:validation}
%=============================================================

This is the reason the OSSE justifies its cost.

The observational estimate of preferential bias post-stratifies by eddy class using altimetry:
with $f_\kappa$ the areal--temporal fraction of class $\kappa$ and $\hat{f}_\kappa$ the
fraction of profiles falling in it,
\begin{equation}
  \label{eq:poststrat}
  \widehat{\mathcal{P}} = \sum_{\kappa} \big(\hat{f}_\kappa - f_\kappa\big)\,\bar\varepsilon_\kappa,
  \qquad \kappa \in \{\text{anticyclone},\ \text{cyclone},\ \text{exterior}\}.
\end{equation}
Applied to the real array, this estimator has never been checked against a known answer.

\paragraph{Run it on the synthetic floats.} The model supplies its own sea surface height field,
so the same eddy detection and tracking algorithm can be applied to model output, producing the
model's analogue of an altimetric eddy atlas. Compute \eqref{eq:poststrat} for the synthetic
profiles and compare against $\widehat{\mathcal{P}}^{\text{sim}}$, which is known. Three
outcomes, all informative:

\begin{itemize}[leftmargin=*]
  \item \textbf{Recovery.} \eqref{eq:poststrat} reproduces the known bias within its error bar
        $\Rightarrow$ the observational estimator is validated and can be applied to real floats
        with a defensible uncertainty.
  \item \textbf{Underestimation.} \eqref{eq:poststrat} is systematically too small
        $\Rightarrow$ preferential structure operates below the resolution of the eddy partition
        (selection by amplitude, by radius, or at submesoscale). The ratio is a correction
        factor; refine the partition and repeat.
  \item \textbf{Failure.} No relationship $\Rightarrow$ the class partition does not capture the
        mechanism, and the observational route must be rethought.
\end{itemize}

\paragraph{Refinement diagnostic.} Recompute \eqref{eq:poststrat} under progressively finer
partitions---adding strata for eddy amplitude quantile and normalised radius from the core.
Saturation indicates the relevant structure has been captured; continued growth bounds how much
\eqref{eq:poststrat} is missing.

\paragraph{Occupancy imbalance as descriptive evidence.} Independently of
\eqref{eq:poststrat}, the comparison of $\hat{f}_\kappa$ against $f_\kappa$ is the most direct
demonstration that preferential sampling occurred at all: under exogenous sampling the two agree
up to sampling error. Report the comparison with confidence intervals, and---per
Remark~\ref{rem:trend}---also year by year, testing the imbalance for trend.

%=============================================================
\section{Internal consistency}
\label{sec:consistency}
%=============================================================

\paragraph{Use the model's own climatology for $\mu$.} Inside the OSSE, $\mu^{\text{model}}$
must be computed from the model's long record. Substituting an external climatology would place
the model-minus-observation offset---a large, static field structured along fronts and
bathymetry, precisely where float behaviour is non-random---into $\varepsilon$, inflating its
variance and correlating it with sampling. That manufactures a preferential signal out of a
model bias.

\paragraph{Compute it the way the real one is computed.} Long-record, per grid cell, by the same
local polynomial regression and with the same resolved/unresolved split---not a shortcut mean
over the analysis window. The OSSE is worth more if it exercises the actual pipeline.

\paragraph{No inconsistency with the real analysis,} where $\mu$ must come from an
Argo-independent long record, per Remark~\ref{rem:circular}. Each experiment is internally
closed: the model's truth is the model, the ocean's truth is the ocean.

\paragraph{Climatology-error sensitivity.} Run the OSSE a second time with a deliberately
imperfect $\mu$---the model's mean from a different period, or an external climatology---to
quantify how much climatology error leaks into the bias estimate. Cheap once the machinery
exists, and a question that will be asked.

%=============================================================
\section{Protocol}
%=============================================================

\begin{description}[leftmargin=*,style=nextline]
  \item[Stage 0 --- Field preparation.] Compute $Y^{\text{model}}$ ($0$--$2000$\,m OHC by
    vertical integration of $\rho c_p T$ over model levels, masked to bathymetry $>2000$\,m) for
    the full period. Compute $\mu^{\text{model}}$ per cell from the long record. Fix the
    resolved space and compute $\varepsilon^{\text{model}}$. Record
    $\sigma = \mathrm{sd}_P(\varepsilon^{\text{model}})$.
  \item[Stage 1 --- Truth.] $\bar{Y}^{\text{sim}} = \int Y^{\text{model}} dP$ by quadrature over
    the masked grid. Note the quadrature error.
  \item[Stage 2 --- Trajectories.] Seed and integrate the ensemble
    (Section~\ref{sec:simulation}). Record every profile's position, time, float ID, and cycle
    number---the last two are required for the effective sample size diagnostic.
  \item[Stage 3 --- Sampling.] Evaluate $Y^{\text{model}}$ and $\varepsilon^{\text{model}}$ at
    synthetic profile positions. Compute $\hat{Y}^{\text{sim}}$ and $b^{\text{sim}}$ per member.
  \item[Stage 4 --- Control.] Evaluate $\varepsilon$ at the same positions from other years
    (Section~\ref{sec:control}). Compute $b^{\text{ctrl}}$ for each $y'$; form
    $\widehat{\mathcal{C}}$, $\widehat{\mathcal{P}}$, and the null distribution.
  \item[Stage 5 --- Nugget structure.] Variograms of $\varepsilon^{\text{model}}$ at synthetic
    profile positions, binned separately for same-float and cross-float pairs. Report
    $\ell_\varepsilon$ and $n_{\text{eff}}$. This also checks the Section~\ref{sec:compatible}
    claim that the nugget has enough spatial structure for trapping to be possible.
  \item[Stage 6 --- Eddy census.] Run the eddy detection algorithm on model sea surface height.
    Compute $f_\kappa$, $\hat{f}_\kappa$, and the post-stratified estimate
    \eqref{eq:poststrat}, both pooled and year by year.
  \item[Stage 7 --- Validation.] Compare \eqref{eq:poststrat} against
    $\widehat{\mathcal{P}}^{\text{sim}}$. Repeat under partition refinement.
  \item[Stage 8 --- Closure and sensitivities.] Budget closure; climatology-error run;
    resolved-space sensitivity; assimilating versus free-running comparison.
\end{description}

%=============================================================
\section{Deliverables}
\label{sec:deliverables}
%=============================================================

\begin{table}[h]
\centering
\begin{tabular}{lcc}
\toprule
Quantity & Free-running & Reanalysis \\
\midrule
$n$, $n_{\text{eff}}$ & & \\
$\sigma$ (J\,m$^{-2}$) & & \\
$\ell_\varepsilon$, same-float vs cross-float & & \\
$\hat{Y} - \bar{Y}$ (total) & & \\
$\widehat{\mathcal{C}}^{\text{sim}}$ & & \\
$\widehat{\mathcal{P}}^{\text{sim}}$ $\pm$ ensemble s.e. & & \\
Control null spread & & \\
$\hat{f}_\kappa$ vs $f_\kappa$ by class & & \\
Occupancy imbalance trend (per year) & & \\
\eqref{eq:poststrat} vs $\widehat{\mathcal{P}}^{\text{sim}}$ & & \\
Budget closure residual & & \\
\bottomrule
\end{tabular}
\caption{Headline table. Run in full for both a quiescent and an energetic box; the contrast is
the generalisable result.}
\end{table}

\paragraph{Figures.}
\begin{enumerate}[leftmargin=*]
  \item Sample synthetic trajectories over model sea surface height---the mechanism, visually.
        Highlight a trapped float and a flushed one.
  \item Distribution of $b^{\text{ctrl}}$ across control years with $b^{\text{sim}}$ marked---the
        test.
  \item Occupancy $\hat{f}_\kappa$ against areal fraction $f_\kappa$ by class, with confidence
        intervals.
  \item \eqref{eq:poststrat} against $\widehat{\mathcal{P}}^{\text{sim}}$ across partition
        refinements---the validation curve.
  \item Residence time in the box against mean $\varepsilon$ encountered, per float---the
        mechanism as a direct association, independent of any null construction.
  \item Occupancy imbalance by year, with trend---the diagnostic of Remark~\ref{rem:trend}.
\end{enumerate}

Figure 5 is worth producing early: it requires no null and no eddy census, and if the
association is absent in the model, that is known before the rest is built.

%=============================================================
\section{Transfer to the real ocean}
\label{sec:transfer}
%=============================================================

What transfers is not $b^{\text{sim}}$ itself but the \emph{relationships}: the validated
estimator \eqref{eq:poststrat} and its correction factor under partition refinement. For those
to carry over, two model properties must be checked against observations.

\paragraph{Nugget variance $\sigma$.} The achievable magnitude of $b$ scales with the
variability of the field being selected on, so a model that understates $\sigma$ understates the
bias. At $1/12^\circ$ with assimilation, GLORYS will. Compare the model's $\varepsilon$ variance
against along-track altimetric variance and against real Argo residuals in the same boxes, and
scale the transfer explicitly.

\paragraph{Lagrangian statistics.} Compare simulated and real distributions of residence time in
the box, per-float profile counts, displacement per cycle, and looping frequency. If synthetic
floats are flushed faster than real ones, trapping---and hence $b$---is understated.

%=============================================================
\section{Limitations}
%=============================================================

\begin{enumerate}[leftmargin=*]
  \item \textbf{Resolution.} $1/12^\circ$ resolves mesoscale but not submesoscale; trapping in
        submesoscale coherent structures is absent entirely. LLC4320 at $1/48^\circ$ partially
        addresses this over a short window.
  \item \textbf{Assimilation damping.} Increments disrupt eddy coherence in the reanalysis,
        shortening Lagrangian trapping times and biasing $b^{\text{sim}}$ toward zero. This is
        the main argument for the free-running run.
  \item \textbf{Float behaviour idealised.} Real floats ground on topography, fail under ice,
        and use varied park-and-profile schedules and parking depths. Not material in North
        Atlantic boxes; material for a global extension, where under-ice failure in the Southern
        Ocean is a strong candidate for the largest preferential bias in the array.
  \item \textbf{Model eddy statistics are not the ocean's.} Eddy amplitude, coherence, and
        polarity asymmetry are model properties, so the simulated occupancy imbalance is
        indicative rather than a measurement of the real array---for which the altimetry route on
        real float positions remains the primary estimate.
  \item \textbf{Split dependence.} $\widehat{\mathcal{C}}$ and $\widehat{\mathcal{P}}$ depend on
        where the resolved/unresolved boundary is drawn; only their sum is invariant. Report
        both, and the sensitivity.
  \item \textbf{Dependence restricted to $\varepsilon$.} Per Section~\ref{sec:compatible} we
        assume the design does not depend on the resolved anomaly $a$. If large-scale warm states
        systematically alter flushing rates, this is violated and the mapped $a$ would itself
        require scrutiny.
\end{enumerate}

%=============================================================
\begin{thebibliography}{99}

\bibitem{chelton2011} Chelton, D.~B., Schlax, M.~G., Samelson, R.~M. (2011).
Global observations of nonlinear mesoscale eddies. \emph{Progress in Oceanography} 91, 167--216.

\bibitem{delandmeter2019} Delandmeter, P., van Sebille, E. (2019).
The Parcels v2.0 Lagrangian framework: new field interpolation schemes.
\emph{Geoscientific Model Development} 12, 3571--3584.

\bibitem{diggle2010} Diggle, P.~J., Menezes, R., Su, T.-L. (2010).
Geostatistical inference under preferential sampling.
\emph{Journal of the Royal Statistical Society C} 59, 191--232.

\bibitem{gelfand2012} Gelfand, A.~E., Sahu, S.~K., Holland, D.~M. (2012).
On the effect of preferential sampling in spatial prediction.
\emph{Environmetrics} 23, 565--578.

\bibitem{kuusela2018} Kuusela, M., Stein, M.~L. (2018).
Locally stationary spatio-temporal interpolation of Argo profiling float data.
\emph{Proceedings of the Royal Society A} 474, 20180400.

\bibitem{lellouche2021} Lellouche, J.-M., et al. (2021).
The Copernicus global $1/12^\circ$ oceanic and sea ice GLORYS12 reanalysis.
\emph{Frontiers in Earth Science} 9, 698876.

\bibitem{park2023} Park, B., Kuusela, M., Giglio, D., Gray, A. (2023).
Spatiotemporal local interpolation of global ocean heat transport using Argo floats:
a debiased latent Gaussian process approach.
\emph{The Annals of Applied Statistics} 17(2), 1491--1520. DOI 10.1214/22-AOAS1679.

\bibitem{rocha2016} Rocha, C.~B., Chereskin, T.~K., Gille, S.~T., Menemenlis, D. (2016).
Mesoscale to submesoscale wavenumber spectra in Drake Passage.
\emph{Journal of Physical Oceanography} 46, 601--620.

\bibitem{rg2009} Roemmich, D., Gilson, J. (2009).
The 2004--2008 mean and annual cycle of temperature, salinity, and steric height in the global
ocean from the Argo Program. \emph{Progress in Oceanography} 82, 81--100.

\bibitem{sukianto2026} Sukianto, T., et al. (2026).
Locally stationary Argo ocean heat content estimates: modeling, validation and uncertainty
quantification. arXiv:2606.31957.

\bibitem{vansebille2018} van Sebille, E., et al. (2018).
Lagrangian ocean analysis: fundamentals and practices. \emph{Ocean Modelling} 121, 49--75.

\end{thebibliography}

\end{document}
